%! Rational Functions and Zeros — Unit 3, Lesson 3.2 — Lesson Plan
\documentclass[10pt]{article}
\usepackage{precalculus-article}
\usepackage{precalculus-boxes}
\usepackage{graphicx}
\graphicspath{{images/}}
\newcommand{\TallMath}[1]{$\displaystyle #1\rule[-1.4em]{0pt}{3.2em}$}
\newcommand{\CourseName}{Precalculus}
\newcommand{\SchoolYear}{2026--2027}
\newcommand{\MeetingLength}{55 minutes}
\newcommand{\UnitNumberName}{Unit 3: Rational Functions \quad}
\newcommand{\LessonNumberName}{Lesson 3.2: Rational Functions and Zeros}

\begin{document}
\begin{center}
    {\Large\bfseries \CourseName: \SchoolYear \\
    {\small\bfseries \UnitNumberName \LessonNumberName}}
\end{center}

\begin{tcolorbox}[colback=sky, colframe=navy, boxrule=0.9pt, arc=2mm,
    left=2mm, right=2mm, top=1.5mm, bottom=1.5mm]
    \textbf{Primary Objective:} Students will identify the real zeros of a rational function
    as the zeros of the numerator that lie in the domain, and use a sign chart to determine
    intervals where a rational function is positive or negative in order to solve rational inequalities.
    \hfill\textit{\small Skills: MP 1.A}
\end{tcolorbox}

\renewcommand{\arraystretch}{1.6}
\begin{skillbox}[Priority Ideas \& Skills]{goldbox}
\begin{tabularx}{\linewidth}{@{}X X@{}}
\textbf{AP Skills} &
\textbf{Key Understandings} \\
\midrule
\textbf{MP 1.A} --- Solve equations and inequalities represented analytically,
with and without technology. &
\begin{itemize}[leftmargin=1.2em,itemsep=2pt,topsep=0pt]
  \item The real zeros of a rational function are the real zeros of the numerator
        that also lie in the domain. (EK 1.8.A.1)
  \item Zeros of the numerator and zeros of the denominator are both critical
        values used as endpoints for sign analysis intervals. (EK 1.8.A.2)
  \item A sign chart organizes critical values into intervals; a test point in
        each interval determines the sign of the function on that interval.
\end{itemize} \\
\end{tabularx}
\end{skillbox}

\begin{skillbox}[{Vocabulary, Concepts, \& Theorems}]{greenbox}
\renewcommand{\arraystretch}{1.4}
\begin{tabularx}{\linewidth}{@{} >{\leavevmode\bfseries\color{navy}}p{0.24\linewidth} X @{}}
Zero of $r(x)$      & A value $x = c$ in the domain where $r(c) = 0$; the real zeros of $r$ are the real zeros of its numerator that are also in the domain. \\
Domain restriction  & A value excluded from the domain because it makes the denominator zero; NOT a zero of the function. \\
Critical value       & Any $x$-value where $r(x) = 0$ (numerator zero in domain) or $r(x)$ is undefined (denominator zero); used as boundaries for sign analysis. \\
Sign chart          & A number-line diagram marking critical values and showing the sign ($+$ or $-$) of $r(x)$ on each resulting interval. \\
Rational inequality & An inequality of the form $r(x) \geq 0$, $r(x) \leq 0$, $r(x) > 0$, or $r(x) < 0$, solved by sign analysis using all critical values. \\
\end{tabularx}
\end{skillbox}

\begin{skillbox}[Activate Prior Knowledge \& Spiral Review]{sky}
    \begin{tabularx}{\linewidth}{@{}X c@{}}
        \begin{minipage}[t]{0.70\linewidth}\vspace{0pt}
            Please see attached warm-up (spiral review).
            Skills reviewed: finding horizontal asymptotes by comparing leading terms (LO 1.7.A),
            interpreting a horizontal asymptote at $y = 0$, factoring a quadratic trinomial.
        \end{minipage}
        &
        \begin{minipage}[t]{0.24\linewidth}\vspace{0pt}\centering
            \textbf{LO 1.7.A}\\Horizontal asymptotes\\[4pt]
            \textbf{LO 1.8.A}\\Zeros of rational functions
        \end{minipage}
    \end{tabularx}
\end{skillbox}

\begin{skillbox}[Hook (3--5 min)]{redbox}
A company models its \textbf{profit per unit} as
\[
  P(x) = \frac{x^2 - 9}{x - 1}, \qquad x > 0,\; x \ne 1,
\]
where $x$ is the number of units produced (in thousands).

\medskip
\textbf{Discussion questions:}
\begin{itemize}[leftmargin=1.4em,itemsep=3pt]
  \item At what production level(s) $x$ does the company break even (profit per unit $= 0$)?
  \item For which values of $x$ is the profit per unit \textit{positive}?
  \item Why can we not simply evaluate $P(1)$?
\end{itemize}
\end{skillbox}

\vspace{0.1in}

\begin{multicols}{2}

\begin{notesbox}{Part 1: Zeros of Rational Functions (10 min)}
\textbf{Key Principle (EK 1.8.A.1).} A rational function
$r(x) = \dfrac{p(x)}{q(x)}$ equals zero when $p(x) = 0$
\textit{and} $q(x) \ne 0$ (i.e., $x$ is in the domain).

\medskip
\textbf{Worked Example.} Find the zeros of
$r(x) = \dfrac{x^2 - 4}{x - 3}$.

\begin{enumerate}[itemsep=4pt]
  \item \textbf{Domain:} $q(x) = x - 3 = 0 \Rightarrow x = 3$.
        Domain is all real numbers except $x = 3$.
  \item \textbf{Set numerator $= 0$:}
        $x^2 - 4 = 0 \Rightarrow (x-2)(x+2) = 0$.
        Solutions: $x = 2$ and $x = -2$.
  \item \textbf{Check domain:} Neither $2$ nor $-2$ equals $3$,
        so both are in the domain.
  \item \textbf{Zeros:} $x = 2$ and $x = -2$.
\end{enumerate}

\medskip
\textit{Note:} $x = 3$ is NOT a zero of $r$; it is a domain restriction.
(Lesson 1.9 will classify it as a vertical asymptote or hole.)
\end{notesbox}

\columnbreak

\begin{notesbox}{Part 2: Sign Chart Method (15 min)}
\textbf{Key Principle (EK 1.8.A.2).} The sign of $r(x)$ can change
only at a zero of the numerator or a zero of the denominator.
Both types of zeros are \textbf{critical values}.

\medskip
\textbf{4-Step Method.} Solve $\dfrac{x - 1}{x + 2} \geq 0$.

\begin{enumerate}[itemsep=4pt]
  \item \textbf{Critical values:}
        Numerator zero: $x = 1$. Denominator zero: $x = -2$.
  \item \textbf{Number line:}
        Mark $x = -2$ (open circle, excluded) and $x = 1$ (closed circle, included).
  \item \textbf{Test each interval:}
        \begin{itemize}[itemsep=2pt]
          \item $x = -3$: $\frac{-4}{-1} = 4 > 0 \Rightarrow (+)$
          \item $x = 0$: $\frac{-1}{2} < 0 \Rightarrow (-)$
          \item $x = 2$: $\frac{1}{4} > 0 \Rightarrow (+)$
        \end{itemize}
  \item \textbf{Solution ($\geq 0$ means $+$ or zero):}\\
        $(-\infty,\,-2)\cup[1,\,\infty)$
\end{enumerate}
\end{notesbox}

\end{multicols}

\begin{notesbox}{Part 3: Connecting EK 1.8.A.2 (5 min)}
Both zeros of the \textit{numerator} and zeros of the \textit{denominator} are
\textbf{critical values} that serve as interval endpoints in sign analysis.

\medskip
\renewcommand{\arraystretch}{1.4}
\begin{tabular}{@{} l l l @{}}
\textbf{Critical value type} & \textbf{Include in solution?} & \textbf{Reason} \\
\midrule
Numerator zero in domain     & Yes, if $\geq$ or $\leq$       & $r(x) = 0$ satisfies $r \geq 0$ \\
Denominator zero (excluded)  & \textbf{Never}                 & $r(x)$ is undefined there \\
\end{tabular}
\end{notesbox}

\vspace{0.1in}

\begin{skillbox}[Active Monitoring]{redbox}
Circulate and look for:
\begin{itemize}[leftmargin=1.4em,itemsep=3pt]
  \item Students who set the \textit{denominator} equal to zero to find ``zeros'' --- redirect to the numerator.
  \item Students who include denominator-zero values in solution sets --- remind those points are excluded.
  \item Students who forget to check whether numerator zeros are actually in the domain.
\end{itemize}
Cold-call: ``Is $x = 3$ a zero of $r(x) = (x-3)/(x-3)$? Why or why not?''
\end{skillbox}

\begin{skillbox}[Timing \& Pacing]{goldbox}
\renewcommand{\arraystretch}{1.4}
\begin{tabularx}{\linewidth}{@{} >{\bfseries}p{0.14\linewidth} p{0.14\linewidth} X @{}}
Time & Duration & Activity \\
\midrule
0:00--0:05  & 5 min  & Warm-up (spiral: asymptotes, factoring) \\
0:05--0:10  & 5 min  & Hook: profit-per-unit discussion \\
0:10--0:20  & 10 min & Part 1: Zeros of rational functions \\
0:20--0:35  & 15 min & Part 2: Sign chart method (notes + guided practice) \\
0:35--0:40  & 5 min  & Part 3: EK 1.8.A.2 summary table \\
0:40--0:50  & 10 min & Tiered group activity \\
0:50--0:55  & 5 min  & Exit ticket \\
\end{tabularx}
\end{skillbox}

\begin{skillbox}[Reinforcement \& Extension]{goldbox}
See attached homework for 5 targeted problems (zeros, sign charts, reasoning about critical values,
AP-style MC). Lesson 1.9 extends this to vertical asymptotes and holes --- preview by asking:
``What happens to the graph of $r$ near a denominator zero?''
\end{skillbox}

\end{document}
